\begin{abstract}
Multiple-choice scores are commonly interpreted against chance. We argue that this comparison is often insufficient: before an arm is compared with another arm, its score should be tested against the construct's best constant, input-blind predictor. This floor is a necessary, not sufficient, validity condition. Across ten target constructs and one negative control, the floor differs from chance and, for content scoring, is under-specified by the tie convention, length unit, and tokenizer; when option count varies, even ``chance'' is ambiguous. On MMLU-Pro, all 21 evaluated cells are powered at the scale of the reference effect. Among 15 designated damaged cells in four model families, the null choice alone reverses the reading: damaged arms appear above a chance line while only one clears its best-constant floor, and the honest aggregate is 14/15 at or below the floor. Across five smaller benchmarks, 10/15 designated cells show the same wrong-null flip, while a mandatory power analysis explains why most per-benchmark significance tests are inconclusive. We then introduce a complementary, arm-conditional permutation null for a different question: whether one arm's predictions contain item-level information. We prove that its statistic is exactly zero for every pure constant emitter, independent of collapse letter, and use it to re-judge 27 existing cells. The re-analysis withdraws one above-floor capability label, dissolves both below-floor competence labels, and exposes a weak intact-family anchor. Finally, full-fp32 evaluation removes every bf16 exact tie and changes 18.03\% of letter decisions without improving accuracy, ruling out a numerical-tie explanation for the measurement failure. Together, the results motivate a simple reporting rule: calibrate each reported construct to its explicit input-blind null before interpreting arm differences.
\end{abstract}
